\section{Véc-tơ trong không gian}
\subsection{Đề 1}
\Opensolutionfile{ans}[ans/ans-2-B6]
\caulc
%%%%%----------Câu 1
\begin{ex}%[2H2N1-1]
	Trong không gian, gọi $\varphi$ là góc giữa hai vectơ  $\overrightarrow{a}$ và $\overrightarrow{b}$ khác vectơ không. Khẳng định nào sau đây đúng?
	\choice
	{$0^\circ <\varphi<90^\circ$}
	{\True $0^\circ \le \varphi\le 180^\circ$}
	{$0^\circ <\varphi<180^\circ$}
	{$0^\circ \le\varphi\le 90^\circ$}
	\loigiai{
		Gọi $\varphi$ là góc giữa hai vectơ $\overrightarrow{a}$ và $\overrightarrow{b}$. Khi đó $0^\circ \le \varphi\le 180^\circ$.
	}	
\end{ex}
%%%%%----------Câu 2
\begin{ex}%[2H2N1-1]
	Trong các mệnh đề sau, mệnh đề nào \textbf{sai}?
	\choice
	{Nếu hai vectơ vuông góc với nhau thì tích vô hướng của chúng bằng $0$}
	{Tích vô hướng của hai vectơ bằng tích độ dài của hai vectơ đó với côsin góc hợp bởi hai vectơ đó}
	{\True Tích vô hướng của hai vectơ bằng bình phương độ dài của mỗi vectơ}
	{Bình phương vô hướng bằng bình phương độ dài}
	\loigiai{
		Tích vô hướng của hai véc-tơ $\overrightarrow{a}$ và $\overrightarrow{b}$ tính theo công thức $\overrightarrow{a}\cdot \overrightarrow{b}=|\overrightarrow{a}|\cdot |\overrightarrow{b}|\cdot \cos\left(\overrightarrow{a},\overrightarrow{b}\right)$.
		\begin{itemize}[\color{blue}\checkmark]
			\item Nếu $\vec{a}\perp\vec{b}$ thì $\heva{& \left(\vec{a},\vec{b}\right)=90^\circ\\ & \cos\left(\overrightarrow{a},\overrightarrow{b}\right)=0}$. 
			Khi đó $\vec{a}\perp\vec{b} \Leftrightarrow \vec{a}\cdot\vec{b}=0$.
			\item Nếu $\vec{a}=\vec{b}$ thì $\heva{& \left(\vec{a},\vec{b}\right)=0^\circ\\ & \cos\left(\overrightarrow{a},\overrightarrow{b}\right)=1}$. 
			Khi đó $\vec{a}\cdot\vec{b} = \vec{a}\cdot\vec{a}=\vec{a}^2 = |\vec{a}|^2$.
		\end{itemize}
	}
\end{ex}
%%%%%----------Câu 3
\begin{ex}%[2H2N1-1]
	Cho hình hộp $ ABCD.A'B'C'D'$. Đẳng thức nào sau đây là \textbf{sai}?
	\choice
	{$\overrightarrow{AB}+\overrightarrow{AD}+\overrightarrow{AA'}=\overrightarrow{AC'}$}
	{$\overrightarrow{BC}+\overrightarrow{CD}+\overrightarrow{BB'}=\overrightarrow{BD'}$}
	{$\overrightarrow{CB}+\overrightarrow{CD}+\overrightarrow{DD'}=\overrightarrow{CA'}$}
	{\True $\overrightarrow{AD}+\overrightarrow{AB}+\overrightarrow{AA'}=\overrightarrow{A'C}$}
	\loigiai{
		\begin{center}
			\begin{tikzpicture}[scale=1,>=stealth, font=\footnotesize, line join=round, line cap=round, thick, blue]
				\def\a{3}
				\def\h{3}
				\path 	(0:0) coordinate (A)
				++(0:\a) coordinate (D)
				++(-130:\a/2) coordinate (C)
				($(A)+(C)-(D)$) coordinate (B)	
				($(A)!0.25!(C)$) coordinate (H)
				($(A)+(80:\h)$) coordinate (A')
				($(A')+(C)-(A)$) coordinate (C')
				($(C')+(B)-(C)$) coordinate (B')
				($(A')+(D)-(A)$) coordinate (D');
				\draw[dashed] 	(B)--(A)--(D)	(A)--(A')  (A')--(C);
				\draw[] 	(C)--(C') 	(D)--(D') 	(B)--(B')  (B)--(C)--(D) (A')--(B')--(C')--(D')--cycle;
				\foreach \x/\g in {A/180,B/180,C/0,D/0,A'/180,B'/180,C'/0,D'/0}
				\fill[black] (\x) circle (1pt) ($(\g:3mm)+(\x)$) node {$\x$};
			\end{tikzpicture}
		\end{center}
		\begin{itemchoice}
			\itemch Đẳng thức $\vec{AB}+\vec{AD}+\vec{AA'}=\vec{AC'}$\\
			$\vec{AB}+\vec{AD}+\vec{AA'}=(\vec{AB}+\vec{AD})+\vec{AA'} =\vec{AC}+\vec{AA'} =\vec{AC'}$
			\itemch Đẳng thức $\vec{BC}+\vec{CD}+\vec{BB'}=\vec{BD'}$\\
			$\vec{BC}+\vec{CD}+\vec{BB'}=(\vec{BC}+\vec{CD})+\vec{BB'} =\vec{BD}+\vec{BB'} =\vec{BD'}$
			\itemch Đẳng thức $\vec{CB}+\vec{CD}+\vec{DD'}=\vec{CA'}$\\
			$\vec{CB}+\vec{CD}+\vec{DD'}=(\vec{CB}+\vec{CD})+\vec{DD'} =\vec{CA}+\vec{DD'} =\vec{CA'}$
			\itemch Đẳng thức $\vec{AD}+\vec{AB}+\vec{AA'}=\vec{A'C}$\\
			$\vec{AB}+\vec{AD}+\vec{AA'}=(\vec{AB}+\vec{AD})+\vec{AA'} =\vec{AC}+\vec{AA'} =\vec{AC'}\ne\vec{A'C}$
		\end{itemchoice}
	}
\end{ex}
%%%%%----------Câu 4
\begin{ex}%[2H2H1-2]
	Cho hình chóp $S.ABCD$ có đáy $ABCD$ là hình bình hành. Đặt $\overrightarrow{SA}=\overrightarrow{a}$, $\overrightarrow{SB}=\overrightarrow{b}$, 
	$\overrightarrow{SC}=\overrightarrow{c}$, $\overrightarrow{SD}=\overrightarrow{d}$. Khẳng định nào sau đây đúng?
	\choice
	{\True
		$\overrightarrow{a}+\overrightarrow{c}= \overrightarrow{d}+\overrightarrow{b}$}
	{$\overrightarrow{a}+\overrightarrow{b}= \overrightarrow{c}+\overrightarrow{d}$}
	{$\overrightarrow{a}+\overrightarrow{d}= \overrightarrow{b}+\overrightarrow{c}$}
	{$\overrightarrow{a}+\overrightarrow{c} + \overrightarrow{d}+\overrightarrow{b} =\overrightarrow{0}$}
	\loigiai{
		\begin{center}	
			\begin{tikzpicture}[scale=1,>=stealth, font=\footnotesize, line join=round, line cap=round, thick, blue]
				\def\a{1}
				\path (0:0) coordinate (A)
				++(0:4*\a) coordinate (D)
				++(210:2*\a) coordinate (C)
				($(A)+(C)-(D)$) coordinate (B)
				($(A)!.5!(C)$) coordinate (O)
				($(O)+(90:4.25*\a)$) coordinate (S)
				(intersection of A--C and B--D) coordinate (O);%giao điểm O
				\draw[dashed] (B)--(A)--(D)--cycle (A)--(S)--(O) (A)--(C);
				\draw[] (B)--(C)--(D) (B)--(S) (C)--(S) (D)--(S);
				\foreach \x/\y in {A/135,B/-135,C/-45,D/45,S/90,O/-120}
				\fill[black] (\x) circle (1pt) ($(\y:3mm)+(\x)$) node {$\x$};
				\path 
				(S)--(A) node[above,pos=.55,sloped]{$\vec{a}$}
				(S)--(B) node[above,pos=.5,sloped]{$\vec{b}$}
				(S)--(C) node[above,pos=.5,sloped]{$\vec{c}$}
				(S)--(D) node[above,pos=.5,sloped]{$\vec{d}$}
				(S)--(O) node[above,pos=.625,sloped,magenta]{\scriptsize$\vec{a}+\vec{c}=\vec{b}+\vec{d}$}
				;	
			\end{tikzpicture}
		\end{center}
		Gọi $O$ là tâm hình bình hành $ABCD$, khi đó 
		$\overrightarrow{SA}+\overrightarrow{SC}= \overrightarrow{SD}+\overrightarrow{SB} =2\overrightarrow{SO}$.\\
		Vậy $\overrightarrow{a}+\overrightarrow{c}= \overrightarrow{d}+\overrightarrow{b}$
	}
\end{ex}
%%%%%----------Câu 5
\begin{ex}%[2H2H1-2]
	Cho tứ diện $ABCD$. Gọi $M$, $N$, $P$, $Q$, $I$, $J$ lần lượt là trung điểm của $AB$, $BC$, $CD$, $DA$, $AC$, $BD$. Các vectơ bằng nhau là
	\choice
	{$\overrightarrow{MI}$, $\overrightarrow{IQ}$, $\overrightarrow{QM}$}
	{$\overrightarrow{MN}$, $\overrightarrow{CI}$, $\overrightarrow{QP}$}
	{\True $\overrightarrow{MQ}$, $\overrightarrow{NP}$, $\dfrac{1}{2}\left( \overrightarrow{CD}- \overrightarrow{CB} \right)$}
	{$\overrightarrow{MQ}$, $\overrightarrow{NP}$, $\dfrac{1}{2}\left( \overrightarrow{CB}- \overrightarrow{CD} \right)$}
	\loigiai{
		\immini{
			\begin{itemize}
				\item $\overrightarrow{MI} \ne \overrightarrow{IQ}$.
				\item $\overrightarrow{MN} = -\overrightarrow{CI} \ne\vec{CI}$.
				\item $\overrightarrow{MQ}= \overrightarrow{NP} =\dfrac{1}{2}\left(\overrightarrow{CD}-\overrightarrow{CB}\right) =\dfrac{\overrightarrow{BD}}{2}$.
				\item $\overrightarrow{MQ }=\overrightarrow{NP} \ne\dfrac{1}{2}\left( \overrightarrow{CB}- \overrightarrow{CD} \right) =\dfrac{1}{2}\vec{DB}$.
			\end{itemize}
		}
		{	
			\begin{tikzpicture}[scale=1,>=stealth, font=\footnotesize, line join=round, line cap=round, thick,>=stealth]
				\def\a{4}
				\path 	(0:0) coordinate (B)
				++(0:\a) coordinate (D)
				++(-120:\a/2) coordinate (C)
				($(B)+(70:\a)$) coordinate (A)
				($(A)!.5!(B)$) coordinate (M)
				($(B)!.5!(C)$) coordinate (N)
				($(C)!.5!(D)$) coordinate (P)
				($(D)!.5!(A)$) coordinate (Q)
				($(A)!.5!(C)$) coordinate (I)
				($(B)!.5!(D)$) coordinate (J)
				;
				\draw[dashed,->] (J)--(D) (B)--(J);
				\draw[] (B)--(A)--(D) (A)--(C)
				(B)--(C)--(D);
				\draw[magenta,dashed,->] (M)--(Q);
				\draw[magenta,dashed,->] (N)--(P);
				\foreach \x/\y in {A/90,B/180,C/-45,D/0, M/135,N/-120,P/-30,Q/30, I/30,J/120}
				\fill[black] (\x) circle (1pt) ($(\y:3mm)+(\x)$) node {$\x$};
			\end{tikzpicture}
		}
	}
\end{ex}
%%%%%----------Câu 6
\begin{ex}%[2H2H1-2]
	Cho hình lăng trụ $ABC.A'B'C'$. Đặt $\overrightarrow{AA'}=\overrightarrow{a}$, $\overrightarrow{AB}=\overrightarrow{b}$, $\overrightarrow{AC}=\overrightarrow{c}$, 
	$\overrightarrow{BC}=\overrightarrow{d}$. Trong các đẳng thức sau đẳng thức nào đúng?
	\choice
	{$\overrightarrow{a} +\overrightarrow{b}+\overrightarrow{c}+\overrightarrow{d}=\overrightarrow{0}$}
	{$\overrightarrow{a} +\overrightarrow{b}+\overrightarrow{c}=\overrightarrow{d}$}
	{\True $\overrightarrow{b}-\overrightarrow{c}+\overrightarrow{d}=\overrightarrow{0} $}
	{$\overrightarrow{a} =\overrightarrow{b}+\overrightarrow{c}$}
	\loigiai{
		\begin{center}
			\begin{tikzpicture}[line cap=round, line join=round, font=\footnotesize, thick, blue]
				\def\a{4} \def\h{4.5}
				\path 	(0:0) coordinate (A)
						++(0:\a) coordinate (C)
						++(-150:3*\a/4) coordinate (B)
						($(A)+(80:\h)$) coordinate (A')
						($(B)+(80:\h)$) coordinate (B')
						($(C)+(80:\h)$) coordinate (C');
				\draw[dashed] (A)--(C);
				\draw[] (C)--(C') (B)--(B') (A)--(A') (A)--(B)--(C) (A')--(B')--(C')--cycle;
				\foreach \x/\y in {A/180,B/-45,C/0,A'/180,B'/-45,C'/0}
				\fill[black] (\x) circle (1pt) ($(\y:3mm)+(\x)$) node {$\x$};
				\path 
				(A)--(A') node[above,pos=.5,sloped]{$\vec{a}$}
				(A)--(B) node[below,pos=.5,sloped]{$\vec{b}$}
				(A)--(C) node[below,pos=.55,sloped]{$\vec{c}$}
				(B)--(C) node[below,pos=.5,sloped]{$\vec{d}$}
				;	
			\end{tikzpicture}
		\end{center}
		Ta có $\overrightarrow{b}-\overrightarrow{c}+\overrightarrow{d} 
		= \overrightarrow{AB}-\overrightarrow{AC}+\overrightarrow{BC}=
		\overrightarrow{CB}+\overrightarrow{BC}= \overrightarrow{0}$.
	}
\end{ex}
%%%%%----------Câu 7
\begin{ex}%[2H2H1-2]
	Cho hình lập phương $ABCD.A'B'C'D'$ có cạnh bằng $a$. Tính tích $\overrightarrow{AC}\cdot\overrightarrow{B'C'}$.
	\choice
	{$\overrightarrow{AC}\cdot\overrightarrow{B'C'}=\dfrac{\sqrt{2}}{2}a^2$}
	{\True $\overrightarrow{AC}\cdot\overrightarrow{B'C'}=a^2$}
	{$\overrightarrow{AC}\cdot\overrightarrow{B'C'}=\sqrt{2}a^2$}
	{$\overrightarrow{AC}\cdot\overrightarrow{B'C'}=2a^2$}
	\loigiai{
		\immini{
			Do $\vec{AD}=\vec{B'C'}$,\\ 
			nên $\left(\vec{AC},\vec{B'C'}\right) =\left(\vec{AC},\vec{AD}\right) =\widehat{CAD}=45^{\circ}$.\\
			Cho nên
			\begin{align*}
				\vec{AC}\cdot\vec{B'C'}&=\left|\vec{AC}\right|\cdot\left|\vec{B'C'}\right|\cdot\cos \left(\vec{AC},\vec{B'C'}\right)\\
				&=a\sqrt{2}\cdot a\cdot\cos 45^{\circ}=a^2.
			\end{align*}
		}
		{
			\begin{tikzpicture}[scale=1,>=stealth, font=\footnotesize, line join=round, line cap=round, thick, blue]
				\def\a{3}
				\path 	(0:0) coordinate (A)
				++(0:\a) coordinate (D)
				++(-130:\a/2) coordinate (C)
				($(A)+(C)-(D)$) coordinate (B)
				($(A)+(90:\a)$) coordinate (A')
				($(B)+(90:\a)$) coordinate (B')
				($(C)+(90:\a)$) coordinate (C')
				($(D)+(90:\a)$) coordinate (D');
				\draw[dashed] (B)--(A)--(D) (A)--(A') (A)--(C);
				\draw[] (C)--(C') (D)--(D') (B)--(B') (B)--(C)--(D) (A')--(B')--(C')--(D')--cycle ;
				\foreach \x/\y in {A/180,B/180,C/0,D/0,A'/180,B'/180,C'/0,D'/0}
				\fill[black] (\x) circle (1pt) ($(\y:3mm)+(\x)$) node {$\x$};	
			\end{tikzpicture}
		}
	}
\end{ex}
%%%%%----------Câu 8
\begin{ex}%[2H2H1-2]
	Cho hình chóp $S.ABCD$ có đáy $ABCD$ là hình bình hành tâm $O$. Hãy chỉ ra đẳng thức \textbf{sai} trong các đẳng thức sau
	\choice
	{$\vec{SA}+\vec{SC}=2\vec{SO}$}
	{$\vec{SB}+\vec{SD}=2\vec{SO}$}
	{\True $\vec{SA}+\vec{SB}+\vec{SC}+\vec{SD}=\vec{AC}+\vec{BD}$}
	{$\vec{SA}+\vec{SC}=\vec{SB}+\vec{SD}$}
	\loigiai{
		\begin{center}	
			\begin{tikzpicture}[scale=1,>=stealth, font=\footnotesize, line join=round, line cap=round, thick, blue]
				\def\a{4}
				\path 	(0:0) coordinate (A)
				++(0:\a) coordinate (D)
				++(210:\a/2) coordinate (C)
				($(A)+(C)-(D)$) coordinate (B)
				($(A)+(80:\a)$) coordinate (S)
				($(A)!.5!(C)$) coordinate (O)
				(intersection of A--C and B--D) coordinate (O);%giao điểm O
				\draw[dashed] (B)--(A)--(D)--cycle	(A)--(S)--(O) (A)--(C);
				\draw[] (B)--(C)--(D) (B)--(S) (C)--(S)	(D)--(S);
				\foreach \x/\y in {A/135,B/-135,C/-45,D/45,S/90,O/-120}
				\fill[black] (\x) circle (1pt) ($(\y:3mm)+(\x)$) node {$\x$};	
			\end{tikzpicture}
		\end{center}
		Ta có $\vec{SA}+\vec{SC}=2\vec{SO}=\vec{SB}+\vec{SD}$.
		\begin{itemchoice}
			\itemch $\vec{SA}+\vec{SC}=2\vec{SO}$\\
			Vì $O$ là trung điểm $AC$ nên $\forall\text{ điểm } S, \vec{SA}+\vec{SC}=2\vec{SO}$.
			\itemch {$\vec{SB}+\vec{SD}=2\vec{SO}$}\\
			Vì $O$ là trung điểm $BD$ nên $\forall\text{ điểm } S, \vec{SB}+\vec{SD}=2\vec{SO}$.
			\itemch {$\vec{SA}+\vec{SB}+\vec{SC}+\vec{SD}=\vec{AC}+\vec{BD}$}\\
			$\vec{SA}+\vec{SB}+\vec{SC}+\vec{SD}=4\vec{SO}\ne\vec{AC}+\vec{BD}$
			\itemch {$\vec{SA}+\vec{SC}=\vec{SB}+\vec{SD}$}\\
			Vì $O$ là trung điểm $AC$ và $BD$ nên $\forall\text{ điểm } S, \vec{SA}+\vec{SC}=\vec{SB}+\vec{SD}=2\vec{SO}$.
		\end{itemchoice}
	}
\end{ex}
%%%%%----------Câu 9
\begin{ex}%[2H2H1-2]
	Gọi $G$ là trọng tâm của tứ diện $ABCD$. Trong các khẳng định sau, khẳng định nào \textbf{sai}?
	\choice
	{$\overrightarrow{AG}=\dfrac{1}{4}\left(\overrightarrow{AB}+\overrightarrow{AC}+\overrightarrow{AD}\right)$}
	{\True $\overrightarrow{AG}=\dfrac{2}{3}\left(\overrightarrow{AB}+\overrightarrow{AC}+\overrightarrow{AD}\right)$}
	{$\overrightarrow{GA}+\overrightarrow{GB}+\overrightarrow{GC}+\overrightarrow{GD}=\overrightarrow{0}$}
	{$\overrightarrow{OG}=\dfrac{1}{4}\left(\overrightarrow{OA}+\overrightarrow{OB}+\overrightarrow{OC}+\overrightarrow{OD}\right)$}
	\loigiai{
		\begin{center}
			\begin{tikzpicture}[line cap=round, line join=round, thick, blue, font=\footnotesize, >=stealth]
				\def\a{1}
				\path 
				(0:0) coordinate (B)
				(0:4*\a) coordinate (D)
				(310:2*\a) coordinate (C)
				($(B)+(80:3.25*\a)$) coordinate (A)
				(barycentric cs:A=1,B=1,C=1,D=1) coordinate (G)
				;
				\draw[dashed,thick] (B)--(D);
				\draw[] (B)--(A)--(D) (A)--(C) (B)--(C)--(D);
				\foreach \x/\y in {A/90,B/180,C/-45,D/0,G/45}
				\fill[black] (\x) circle (1pt) ($(\y:3mm)+(\x)$) node {$\x$};
				\foreach \x/\y in {G/A,G/B,G/C,G/D}{
				\draw[magenta,dashed,->] (\x)--(\y);}	
			\end{tikzpicture}
		\end{center}
		Ta có
		\begin{itemchoice}
			\itemch {$\vec{AG} =\dfrac{1}{4}\left(\vec{AB}+\vec{AC}+\vec{AD}\right)$}\\
			Ta có $\begin{aligned}[t]
				& \vec{GA}+\vec{GB}+\vec{GC}+\vec{GD}=\vec{0}\\
				\Leftrightarrow &\ \vec{AG}=\vec{GB}+\vec{GC}+\vec{GD}\\
				\Leftrightarrow &\ \vec{AG}=\left(\vec{AB}-\vec{AG}\right)+\left(\vec{AC}-\vec{AG}\right)+\left(\vec{AD}-\vec{AG}\right)\\
				\Leftrightarrow &\ \vec{AG}=\dfrac{1}{4}\left(\vec{AB}+\vec{AC}+\vec{AD}\right).
			\end{aligned}$
			\itemch { $\vec{AG} =\dfrac{2}{3}\left(\vec{AB}+\vec{AC}+\vec{AD}\right)$}\\
			Sai khác.
			\itemch {$\vec{GA}+\vec{GB}+\vec{GC}+\vec{GD}=\vec{0}$}\\
			Luôn luôn đúng.
			\itemch {$\vec{OG} =\dfrac{1}{4}\left(\vec{OA}+\vec{OB}+\vec{OC}+\vec{OD}\right)$}\\
			Điểm $G$ là trọng tâm của tứ diện $ABCD$ ta luôn có $\vec{GA}+\vec{GB}+\vec{GC}+\vec{GD}=\vec{0}$.\\
			Khi đó $\begin{aligned}[t]
				& \vec{GA}+\vec{GB}+\vec{GC}+\vec{GD}=\vec{0}\\
				\Leftrightarrow &\ \vec{OA}-\vec{OG} +\vec{OB}-\vec{OG} +\vec{OC}-\vec{OG} +\vec{OD}-\vec{OG}=\vec{0}\\
				\Leftrightarrow &\ 4\vec{OG} =\vec{OA}+\vec{OB}+\vec{OC}+\vec{OD}\\
				\Leftrightarrow &\ \vec{OG} =\dfrac{1}{4}\left(\vec{OA}+\vec{OB}+\vec{OC}+\vec{OD}\right).
			\end{aligned}$
		\end{itemchoice}
	}
\end{ex}
%%%%%----------Câu 10
\begin{ex}%[2H2H1-3]
	Cho hình chóp $S.ABCD$ có đáy $ABCD$ là hình bình hành và $AB=6$, $AD=4$, $\overrightarrow{AB}\cdot\overrightarrow{AD}=-12$ . Tính $\left(\overrightarrow{SC} -\overrightarrow{SA}\right)^2$.
	\choice
	{$76$}
	{\True$28$}
	{$52$}
	{$40$}
	\loigiai{
		Ta có $\begin{aligned}[t]
			\left(\vec{SC}-\vec{SA}\right)^2 = \vec{AC}^2 &= \left(\vec{AB} -\vec{BC}\right)^2\\
			&= \vec{AB}^2+\vec{AD}^2+2 \vec{AB}\cdot\vec{AD}\\
			&= 36+16-2=28.
		\end{aligned}$
	}
\end{ex}
%%%%%----------Câu 11
\begin{ex}%[2H2H1-3]
	Cho tứ diện $ABCD$ có tam giác $ BCD$ đều và $AD=AC$. Giá trị của $\cos\left(\overrightarrow{AB}, \overrightarrow{CD} \right)$ là
	\choice
	{$\dfrac{1}{2}$}
	{\True$0$}
	{$-\dfrac{1}{2}$}
	{$\dfrac{\sqrt{3}}{2}$}
	\loigiai{
		\immini{
			Gọi $M$ là trung điểm của $CD$.\\
			Ta có
			$\triangle BCD$ đều nên $BM\perp CD$, $\triangle ACD$ cân tại $A$ nên $AM\perp CD$.\\
			Suy ra \\	$\vec{AB}\cdot \vec{CD} =\left(\vec{AM}+\vec{MB} \right)\cdot\vec{CD} =\vec{AM}\cdot \vec{CD} +\vec{MB}\cdot \vec{CD}=0$.\\
			$\Rightarrow\cos\left(\vec{AB}, \vec{CD} \right) =\dfrac{\vec{AB}\cdot\vec{CD}}{\left|\vec{AB}\right|\cdot\left|\vec{CD} \right|} =0$
		}{
			\begin{tikzpicture}[scale=1,>=stealth, font=\footnotesize, line join=round, line cap=round, thick, blue]
				\def\a{4}
				\path 	(0:0) coordinate (B)
				++(0:\a) coordinate (D)
				++(-130:\a/2) coordinate (C)
				($(B)+(70:\a)$) coordinate (A)
				($(C)!.5!(D)$) coordinate (M)
				;
				\draw[dashed] (B)--(D) (B)--(M);
				\draw[] (B)--(A)--(D) (A)--(C) (B)--(C)--(D) (A)--(M);
				\foreach \x/\y in {A/90,B/180,C/-45,D/0,M/-30}
				\fill[black] (\x) circle (1pt) ($(\y:3mm)+(\x)$) node {$\x$};
			\end{tikzpicture}
		}	
	}
\end{ex}
%%%%%----------Câu 12
\begin{ex}%[2H2H1-2]
	Cho tứ diện $ABCD$. Gọi $M$ , $N$ lần lượt là trung điểm của $AB$, $CD$. Biết $\overrightarrow{MN} = k\left(\overrightarrow{AD} + \overrightarrow{BC} \right)$. Giá trị $k$ là
	\choice
	{$3$}
	{$2$}
	{$\dfrac{1}{3}$}
	{\True $\dfrac{1}{2}$}
	\loigiai{
		Vì $M$, $N$ lần lượt là trung điểm của $AB$, $CD$  nên ta có:\\ $\vec{MA}+\vec{MB}=\vec{0}$, $\vec{DN}+\vec{CN}=\vec{0}$.\\
		Mặt khác $ \heva{
			&\vec{MN} =\vec{MA}+\vec{AD}+\vec{DN}\\
			&\vec{MN} =\vec{MB}+\vec{BC}+\vec{CN}. }$\\
		$\Rightarrow 2\vec{MN} = \vec{AD}+\vec{BC}
		\Rightarrow \vec{MN} = \dfrac{1}{2}\left(\vec{AD} + \vec{BC} \right)$.\\ 
		Vậy $k=\dfrac{1}{2}$
	}
\end{ex}
\Closesolutionfile{ans}
\inputansbox[3]{6}{ans/ans-2-B6}
\cauds
\Opensolutionfile{ans}[ans/ans-2-B6-DS]
\setcounter{ex}{0}
%%%%%----------Câu 13
\begin{ex}%[2H2N1-1]
	Trong không gian, cho ba vec-tơ $\overrightarrow{a}$, $\overrightarrow{b}$, $\overrightarrow{c}$ phân biệt đều khác $\overrightarrow{0}$. 
	\choiceTF[t]
	{\True  Nếu $\overrightarrow{a}$ và $\overrightarrow{b}$ cùng hướng với $\overrightarrow{c}$ thì $\overrightarrow{a}$ và $\overrightarrow{b}$ cùng hướng}
	{Nếu $\overrightarrow{a}$ và $\overrightarrow{b}$ ngược hướng với $\overrightarrow{c}$ thì $\overrightarrow{a}$ và $\overrightarrow{b}$ ngược hướng}
	{Nếu $\overrightarrow{a}$ và $\overrightarrow{b}$ cùng phương với $\overrightarrow{c}$ thì $\overrightarrow{a}$ và $\overrightarrow{b}$ cùng hướng}
	{Nếu  $\overrightarrow{a}$ vuông góc với  $\overrightarrow{b}$ và $\overrightarrow{b}$ vuông góc với  $\overrightarrow{c}$ thì  $\overrightarrow{a}$ vuông góc với  $\overrightarrow{c}$}
	\loigiai{
		\begin{itemchoice}
			\itemch \textbf{Đúng}.
			\begin{center}
				\begin{tikzpicture}[scale=1,>=stealth, font=\footnotesize, line join=round, line cap=round,thick, blue]
					\draw[gray!50,thin] (-4,0)--(4,0) (-4,-1)--(4,-1) (-4,-2)--(4,-2);
					\draw[->,teal] (0,0)-- node [above]{$\overrightarrow{c}$}(3,0);
					\draw[->,magenta] (-2,-1)--node [above]{$\overrightarrow{b}$}(0,-1);
					\draw[->,magenta] (-3,-2)--node [above]{$\overrightarrow{a}$}(2,-2);
				\end{tikzpicture}
			\end{center}
			Ta có $\overrightarrow{a}$ và $\overrightarrow{b}$ cùng hướng với $\overrightarrow{c}$ thì $\overrightarrow{a}$ và $\overrightarrow{b}$ cùng hướng.
			\itemch \textbf{Sai}.
			\begin{center}
				\begin{tikzpicture}[scale=1,>=stealth, font=\footnotesize, line join=round, line cap=round, thick]
					\draw[gray!50,thin] (-4,0)--(4,0) (-4,-1)--(4,-1) (-4,-2)--(4,-2);
					\draw[<-,teal] (0,0)-- node [above]{$\overrightarrow{c}$}(3,0);
					\draw[->] (-2,-1)--node [above]{$\overrightarrow{b}$}(0,-1);
					\draw[->] (-3,-2)--node [above]{$\overrightarrow{a}$}(2,-2);
				\end{tikzpicture}
			\end{center}
			Ta có $\overrightarrow{a}$ và $\overrightarrow{b}$ ngược hướng với $\overrightarrow{c}$ thì $\overrightarrow{a}$ và $\overrightarrow{b}$ cùng hướng.
			\itemch \textbf{Sai}.
			\begin{center}
				\begin{tikzpicture}[scale=1,>=stealth, font=\footnotesize, line join=round, line cap=round,thick]
					\draw[gray!50,thin] (-4,0)--(4,0) (-4,-1)--(4,-1) (-4,-2)--(4,-2);
					\draw[<-,teal] (0,0)-- node [above]{$\overrightarrow{c}$}(3,0);
					\draw[<-] (-2,-1)--node [above]{$\overrightarrow{b}$}(0,-1);
					\draw[->] (-3,-2)--node [above]{$\overrightarrow{a}$}(2,-2);
				\end{tikzpicture}
			\end{center}
			Ta có $\overrightarrow{a}$ cùng phương với  $\overrightarrow{b}$ và $\overrightarrow{b}$ cùng phương với  $\overrightarrow{c}$ thì   $\overrightarrow{a}$ và $\overrightarrow{b}$ có thể cùng hướng hoặc ngược hướng.
			\itemch \textbf{Sai}.
			\immini{
				Không làm mất tính tổng quát. Xét hình chóp $S.ABC$ có đáy $\triangle ABC$ vuông tại $B$ và $SA \perp \left(ABC\right)$ (như hình vẽ bên).\\
				Ta có $\heva{&BC\perp SA\\& BC\perp AB}\Rightarrow BC\perp \left(SAB\right)\Rightarrow BC\perp SB$.\\
				Đặt $\overrightarrow{SA}=\overrightarrow{a}$, $\overrightarrow{BC}=\overrightarrow{b}$,
				$\overrightarrow{SB}=\overrightarrow{c}$.\\
				Khi đó $\overrightarrow{a}\perp \overrightarrow{b}$ và $\overrightarrow{b}\perp \overrightarrow{c}$ nhưng $\overrightarrow{a}\not\perp \overrightarrow{c}$.
			}{
				\begin{tikzpicture}[scale=1,>=stealth, font=\footnotesize, line join=round, line cap=round, thick]
					\def\a{4} \def\h{3}
					\path 	(0:0) coordinate (A)
					++(0:\a) coordinate (C)
					++(-150:4*\a/5) coordinate (B)
					($(A)+(90:\h)$) coordinate (S);
					\draw[] (A)--(B)--node [above]{$\vec{b}$}(C)
					(A)--node [left]{$\vec{a}$}(S)	(B)--node [right]{$\vec{c}$}(S)	(C)--(S);
					\draw[dashed] (A)--(C);
					\foreach \x/\y in {A/180,B/-135,C/-30,S/90}
					\fill[black] (\x) circle (1pt) ($(\x)+(\y:3mm)$) node {$\x$};
					\draw pic[draw,angle radius=2mm]{right angle=B--A--S};
					\draw pic[draw,angle radius=2mm]{right angle=C--B--S};
				\end{tikzpicture}	}
			\end{itemchoice}	}
\end{ex}
%%%%%----------Câu 14
\begin{ex}%[2H2H1-2]
	Trong mặt phẳng $(\alpha)$ cho tứ giác $ABCD$ và một điểm $S$ tùy ý. 
	\choiceTF[t]
	{ $\overrightarrow{AC} + \overrightarrow{BD} = \overrightarrow{AB} + \overrightarrow{CD}$}
	{$\overrightarrow{SA} + \overrightarrow{SC} = \overrightarrow{SB} + \overrightarrow{SD}$}
	{\True Nếu tồn tại điểm $S$ sao cho $\overrightarrow{SA} + \overrightarrow{SC} = \overrightarrow{SB} + \overrightarrow{SD}$ thì $ABCD$ là hình bình hành}
	{$\overrightarrow{OA} + \overrightarrow{OB} + \overrightarrow{OC} + \overrightarrow{OD} = \overrightarrow{0}$ khi và chỉ khi $O$ là giao điểm của $AC$ và $BD$}
	\loigiai{
		\begin{itemchoice}
			\itemch \textbf{Sai}.
			Vì  $\overrightarrow{AC} + \overrightarrow{BD} = \overrightarrow{AB} + \overrightarrow{CD} \Leftrightarrow 
			\overrightarrow{AC} - \overrightarrow{AB} = \overrightarrow{DB} - \overrightarrow{DC} \Leftrightarrow  \overrightarrow{BC}= \overrightarrow{CB}\Leftrightarrow B \equiv C$.
			\itemch \textbf{Sai}.\\
			Gọi $M$ là trung điểm $AC$, $N$ là trung điểm $BD$.\\
			Ta có
			$\overrightarrow{SA} + \overrightarrow{SC}=2\overrightarrow{SM}$;
			$\overrightarrow{SB} + \overrightarrow{SD}=2\overrightarrow{SN}$.\\
			Khi đó $\overrightarrow{SA} + \overrightarrow{SC} = \overrightarrow{SB} + \overrightarrow{SD} \Leftrightarrow \overrightarrow{SM}= \overrightarrow{SN} \Leftrightarrow M\equiv N$.\\
			Điều này không xảy ra với $ABCD$ là một tứ giác bất kì.
			\itemch \textbf{Đúng}.
			Gọi $M$ là trung điểm $AC$, $N$ là trung điểm $BD$.\\
			Ta có
			$\overrightarrow{SA} + \overrightarrow{SC}=2\overrightarrow{SM}$;
			$\overrightarrow{SB} + \overrightarrow{SD}=2\overrightarrow{SN}$.\\
			Khi đó $\overrightarrow{SA} + \overrightarrow{SC} = \overrightarrow{SB} + \overrightarrow{SD} \Leftrightarrow \overrightarrow{SM}= \overrightarrow{SN} \Leftrightarrow M\equiv N$. Khi đó $ABCD$ là hình bình hành.
			\itemch \textbf{Sai}.
			Gọi $M$ là trung điểm $AB$, $N$ là trung điểm $CD$.\\
			$O$ là giao điểm của $AB$ và $CD$. Ta có\\
			$\overrightarrow{OA} + \overrightarrow{OB} + \overrightarrow{OC} + \overrightarrow{OD} = \left(\overrightarrow{OA}+\overrightarrow{OB} \right) 
			+ \left(\overrightarrow{OC}+\overrightarrow{OD}\right) =  2\left(\overrightarrow{OM}+\overrightarrow{ON} \right) =\overrightarrow{0}$.\\
			Điều này không xảy ra với $ABCD$ là một tứ giác bất kì.	
		\end{itemchoice}	}
\end{ex}
%%%%%----------Câu 15
\begin{ex}%[2H2V1-4]
	\immini{
		Một chiếc đèn chùm treo có khối lượng $m=5$ kg được thiết kế với đĩa đèn được giữ bởi bốn đoạn xích $SA$, $SB$, $SC$, $SD$ sao cho $S.ABCD$ là hình chóp tứ giác đều có $\widehat{ASC}=60^\circ$ (tham khảo hình vẽ bên). Trọng lực $\overrightarrow{P}$ tác động lên đèn chùm xác định bởi công thức $\overrightarrow{P}=m\overrightarrow{g}$ trong đó $\overrightarrow{g}$ vec-tơ gia tốc rơi tự do có độ lớn $10\rm\, m/s^2$. 
	}{
		\begin{tikzpicture}[scale=.7,>=stealth, font=\footnotesize, line join=round, line cap=round]
			\tikzset{day/.pic=
				{\draw[shade,bottom color=brown!30,top color= white!30,rounded corners=0.5ex,line width=2.5pt,gray!80]
					(0,0) ellipse ({2.5pt} and {8pt});}
			}
			\def\h{6}
			\def\a{3}
			\def\b{1.5}
			\path
			(0,0) coordinate (O)
			($(O)+(0,\h)$) coordinate (S)
			($(O)+(10:\a cm and \b cm)$)coordinate (M)
			($(O)+(180:\a cm and .8*\b cm)$)coordinate (A)
			($(O)+(0:\a cm and .8*\b cm)$)coordinate (C)
			($(O)+(60:\a cm and .8*\b cm)$)coordinate (B)
			($(O)+(-120:\a cm and .8*\b cm)$)coordinate (D)
			;
			\draw[fill=brown]  (M) arc (10:-190:\a cm and \b cm);
			\draw[fill=white]  (M) arc (10:-190:\a cm and .8*\b cm);
			\draw [fill=white] (M) arc (10:190:\a cm and .8*\b cm);
			\draw[fill,bottom color=black!30,top color= brown!70, ,left color=black!50]  ($(S)-(.5,.3)$) rectangle ($(S)+(.5,.3)$);
			\foreach \m in {0,1,2,...,12}{\pic[rotate=-20] at ($(A)+(.25*\m,.5*\m)$) {day};}
			\foreach \m in {0,1,2,...,14}{\pic[rotate=-15] at ($(D)+(.11*\m,.5*\m)$) {day};}
			\foreach \m in {0,1,2,...,10}{\pic[rotate=5] at ($(B)+(-.15*\m,.5*\m)$) {day};}
			\foreach \m in {0,1,2,...,12}{\pic[rotate=18] at ($(C)+(-.25*\m,.5*\m)$) {day};}
			\foreach \x/\g in {A/180,B/150,C/-45,D/-90,S/90}
			\fill[black] 	(\x) circle (4pt)
			($(\g:7mm)+(\x)$) node {$\x$};
		\end{tikzpicture}	}
	\choiceTF[t]
	{\True Độ lớn của trọng lực $\overrightarrow{P}$ tác động lên chiếc đèn chùm bằng $50$ N}
	{Các vec-tơ $\overrightarrow{SA}$, $\overrightarrow{SB}$, $\overrightarrow{SC}$, $\overrightarrow{SD}$ bằng nhau}
	{Độ lớn của lực căng cho mỗi sợi xích bằng $12{,}5$ N}
	{\True Giữ nguyên khối lượng $m = 5$ kg của đèn, thay đổi cách treo sao cho $\widehat{ASC}=90^\circ$ thì độ lớn của lực căng mỗi sợi xích sẽ lớn hơn}
	\loigiai{
		\begin{itemchoice}
			\itemch \textbf{Đúng}.
			Vì  $\left|\overrightarrow{P}\right|=m\left|\overrightarrow{g}\right|=5\cdot10=50$ N.
			\itemch \textbf{Sai}.
			Vì các vec-tơ $\overrightarrow{SA}$, $\overrightarrow{SB}$, $\overrightarrow{SC}$, $\overrightarrow{SD}$ khác hướng.
			\itemch \textbf{Sai}.	
			\immini{
				Gọi $O$ là giao điểm của $AC$ và $BD$.\\
				Ta có $\triangle ASC$ đều nên $SA= \dfrac{2SO}{\sqrt{3}}$.\\
				Mà $\overrightarrow{P}=\overrightarrow{SA}+\overrightarrow{SB}+\overrightarrow{SC}+\overrightarrow{SD}= 4\overrightarrow{SO}$.\\
				$\Rightarrow \left|\overrightarrow{P} \right|= 4\left|\overrightarrow{SO}\right|=50$.
				$ \Rightarrow SO = \left|\overrightarrow{SO} \right| = 12{,}5$.\\
				$ \Rightarrow \left|\overrightarrow{SA}\right|=$
				$\dfrac{2SO}{\sqrt{3}}= \dfrac{2\cdot 12{,}5}{\sqrt{3}}\approx 14{,}4$.\\ 
				Vậy $\left|\overrightarrow{SA}\right| =\left|\overrightarrow{SB}\right|
				=\left|\overrightarrow{SC}\right|=\left|\overrightarrow{SD}\right|\approx 14{,}4$ N.}
				{\begin{tikzpicture}[scale=.7,>=stealth, font=\footnotesize, line join=round, line cap=round]
					\def\a{4}
					\def\h{4.5}
					\path 	(0:0) coordinate (A)
					++(0:\a) coordinate (D)
					++(-130:\a/2) coordinate (C)
					($(A)+(C)-(D)$) coordinate (B)
					(intersection of A--C and B--D) coordinate (O)
					($(O)+(90:\h)$) coordinate (S);
					\draw[dashed,thick] 	(B)--(A)--(D)	(A)--(S)
					(A)--(C)	(B)--(D)	(S)--(O)	;
					\draw[thick] 			(B)--(C)--(D)
					(B)--(S)	(C)--(S)	(D)--(S);
					\foreach \x/\g in {A/135,B/-135,C/-45,D/45,S/90,O/-90}
					\fill[black] 	(\x) circle (1.5pt)
					($(\g:3mm)+(\x)$) node {$\x$};	
					%\draw pic[draw,angle radius=3mm]{right angle=D--O--S};%Theo chiều dương	
				\end{tikzpicture}
			}
			\itemch Đúng.
			Gọi $O$ là giao điểm của $AC$ và $BD$.\\
			Ta có $\triangle ASC$ vuông cân nên $SA= SO\sqrt{2}$.\\
			Mà $\overrightarrow{P}=\overrightarrow{SA}+\overrightarrow{SB}+\overrightarrow{SC}+\overrightarrow{SD}= 4\overrightarrow{SO}$.\\
			$\Rightarrow \left|\overrightarrow{P} \right|= 4\left|\overrightarrow{SO}\right|=50$.
			$ \Rightarrow SO = \left|\overrightarrow{SO} \right| = 12{,}5$.\\
			$ \Rightarrow \left|\overrightarrow{SA}\right|=$
			$SO\sqrt{2} = 12.5\cdot\sqrt{2}\approx 17{,}7>14{,}4$.\\ 
			Vậy độ lớn của lực căng cho mỗi sợi xích sẽ lớn hơn.\\
		\end{itemchoice}
	}
\end{ex}
%%%%%----------Câu 16
\begin{ex}%[2H2V1-4]
	\immini{
		Cho tứ diện đều $ABCD$. Gọi $M$, $N$, $P$, $Q$, $I$, $J$ lần lượt là trung điểm của $AB$, $BC$, $CD$, $DA$, $AC$, $BD$ (tham khảo hình vẽ bên). 
	}{
		\begin{tikzpicture}[scale=.8,>=stealth, font=\footnotesize, line join=round, line cap=round, thick]
			\def\a{1}
			\path (0:0) coordinate (B)
			(0:5*\a) coordinate (D)
			(300:2.75*\a) coordinate (C)
			($(B)+(80:4*\a)$) coordinate (A)
			($(A)!.5!(B)$) coordinate (M)
			($(B)!.5!(C)$) coordinate (N)
			($(C)!.5!(D)$) coordinate (P)
			($(D)!.5!(A)$) coordinate (Q)
			($(A)!.5!(C)$) coordinate (I)
			($(B)!.5!(D)$) coordinate (J)
			;
			\draw[dashed] (B)--(D) (M)--(Q)--(J)--cycle (N)--(P)--(J)--cycle;
			\draw[] (B)--(A)--(D) (A)--(C) (B)--(C)--(D) (M)--(N)--(I)--cycle (I)--(P)--(Q)--cycle;
			\foreach \x/\y in {A/90,B/180,C/-45,D/30,M/135,N/-150,P/-30,Q/30,I/180,J/50}
			\fill[black] (\x) circle (1pt) ($(\y:3mm)+(\x)$) node {$\x$};
		\end{tikzpicture}	
	}
	\choiceTF[t]
	{ \True Có $3$ vec-tơ bằng $\overrightarrow{MN}$}
	{$\overrightarrow{IM}+\overrightarrow{IN}+\overrightarrow{IP}+\overrightarrow{IQ}=\overrightarrow{IJ}$}
	{\True Hai vec-tơ $\overrightarrow{IM}$ và $\overrightarrow{IP}$ vuông góc với nhau}
	{\True Gọi $O$ là trung điểm của $MP$, $G$ là trọng tâm của tam giác $BCD$. Khi đó $\overrightarrow{AO} = 3 \overrightarrow{OG}$}
	\loigiai{
		\begin{itemchoice}
			\itemch \textbf{Đúng}.
			Các vec-tơ bằng $\overrightarrow{MN}$ là $\overrightarrow{AI}$, $\overrightarrow{IC}$,  $\overrightarrow{QP}$. \\
			Vậy Có $3$ vec-tơ bằng $\overrightarrow{MN}$.
			\itemch \textbf{Sai}.
			\immini{
				Gọi $O$ là giao điểm của $MP$ và $NQ$.\\
				Ta có $\overrightarrow{IM}+\overrightarrow{IN}+\overrightarrow{IP}+\overrightarrow{IQ}=4\overrightarrow{IO}$.\\
				Mà $\overrightarrow{IJ}=2\overrightarrow{IO}$.\\
				Do đó $\overrightarrow{IM}+\overrightarrow{IN}+\overrightarrow{IP}+\overrightarrow{IQ}=2\overrightarrow{IJ}$.
			}{
				\begin{tikzpicture}[>=stealth, font=\footnotesize, line join=round, line cap=round, thick]
					\def\a{.9}
					\def\h{3}
					\path 	(0:0) coordinate (M)
					(0:3*\a) coordinate (Q)
					++(-135:1.5*\a) coordinate (P)
					($(M)+(P)-(Q)$) coordinate (N)
					(intersection of M--P and N--Q) coordinate (O)
					($(O)+(90:\h)$) coordinate (I);
					\draw[dashed] (I)--(M) (N)--(M)--(Q) (M)--(P) (N)--(Q) (I)--(O) ;
					\draw[] (N)--(P)--(Q) (I)--(N) (I)--(P) (I)--(Q);
					\foreach \x/\y in {M/30,N/180,P/0,Q/0,I/90,O/-90}
					\fill[black] (\x) circle (1pt) ($(\y:3mm)+(\x)$) node {$\x$};	
				\end{tikzpicture}
			}	
			\itemch \textbf{Đúng}.
			Ta có $ABCD$ là tứ diện đều nên $BC\perp AD$.\\
			Mà $\overrightarrow{IM}=\dfrac{1}{2}\overrightarrow{CB}$, $\overrightarrow{IP}=\dfrac{1}{2}\overrightarrow{AD}$.\\
			Suy ra $\overrightarrow{IM}$ và $\overrightarrow{IP}$ vuông góc với nhau.
			\itemch \textbf{Đúng}.
			\immini{
				Xét tam giác $ABP$, 
				ta có $\overrightarrow{AM}=\dfrac{1}{2}\overrightarrow{AB}$, 
				$\overrightarrow{PG}=\dfrac{1}{3}\overrightarrow{PB}$.\\
				$\overrightarrow{AO}=\dfrac{1}{2}\left(\overrightarrow{AM}+ \overrightarrow{AP}\right)= \dfrac{1}{4}\overrightarrow{AB}+\dfrac{1}{2}\overrightarrow{AP}$.
				\begin{eqnarray*}
					\overrightarrow{OG}&=& \overrightarrow{OP}+\overrightarrow{PG}\\
					&=&\dfrac{1}{2}\overrightarrow{MP}+\dfrac{1}{3}\overrightarrow{PB}\\
					&=&\dfrac{1}{2}\left(\overrightarrow{AP}-\overrightarrow{AM}\right)+ \dfrac{1}{3}\left(\overrightarrow{AB}-\overrightarrow{AP}\right)\\
					&=& \dfrac{1}{2}\overrightarrow{AP}-\dfrac{1}{4}\overrightarrow{AB}+ \dfrac{1}{3}\overrightarrow{AB}-\dfrac{1}{3}\overrightarrow{AP}\\
					&=& \dfrac{1}{12}\overrightarrow{AB}+ \dfrac{1}{6}\overrightarrow{AP}\\
					&=& \dfrac{1}{3}\left(\dfrac{1}{4}\overrightarrow{AB}+ \dfrac{1}{2}\overrightarrow{AP} \right) = \dfrac{1}{3}\overrightarrow{AO}.
				\end{eqnarray*}
				Vậy $\overrightarrow{AO}=2\overrightarrow{OG}$.
			}{
				\begin{tikzpicture}[scale=.8,>=stealth, font=\footnotesize, line join=round, line cap=round, thick]
					\def\a{4}
					\path (0:0) coordinate (B)
					++(0:\a) coordinate (D)
					++(-150:.9*\a) coordinate (C)
					($(B)+(75:\a)$) coordinate (A)
					($(A)!.5!(B)$) coordinate (M)
					($(B)!.5!(C)$) coordinate (N)
					($(C)!.5!(D)$) coordinate (P)
					($(D)!.5!(A)$) coordinate (Q)
					($(A)!.5!(C)$) coordinate (I)
					($(B)!.5!(D)$) coordinate (J)
					($(M)!.5!(P)$) coordinate (O)
					($(B)!2/3!(P)$) coordinate (G)
					;
					\draw[dashed] (B)--(D) (M)--(P)--(B) (A)--(G);
					\draw[] (B)--(A)--(D) (A)--(C)	(B)--(C)--(D) (A)--(P);
					\foreach \x/\y in {A/90,B/180,C/-45,D/30,M/135,P/-30,O/30,G/-120}
					\fill[black] (\x) circle (1pt) ($(\y:3mm)+(\x)$) node {$\x$};	
				\end{tikzpicture}	}
		\end{itemchoice}	}
\end{ex}
\Closesolutionfile{ans}
\inputansbox[3]{2}{ans/ans-2-B6-DS}
\Opensolutionfile{ans}[ans/ans-2-B6-KQ]
\caukq
\setcounter{ex}{0}
%%%%%----------Câu 17
\begin{ex}%[2H2H1-3]
	Trong không gian, cho hai vec-tơ $\overrightarrow{a}$ và $\overrightarrow{b}$ có cùng độ dài bằng $2$. Biết rằng góc giữa hai vec-tơ đó là $30^\circ$. Tính $\overrightarrow{a} \cdot \overrightarrow{b}$.
	\shortans{$2$}
	\loigiai{
		Ta có $\overrightarrow{a} \cdot \overrightarrow{b}= \left| \overrightarrow{a}\right| \cdot \left| \overrightarrow{b}\right|\cdot
		\cos \left(\overrightarrow{a},\overrightarrow{b} \right) = 2\cdot 2 \cdot \cos 30^\circ = 2$.	}
\end{ex}
%%%%%----------Câu 18
\begin{ex}%[2H2H1-3]
	Cho hình lập phương $ABCD.A'B'C'D'$. Góc giữa hai vec-tơ $\overrightarrow{AC}$ và $\overrightarrow{B'C'}$ bằng bao nhiêu độ? 
	\shortans{$45$}
	\loigiai{
		\begin{center}
			\begin{tikzpicture}[>=stealth, font=\footnotesize, line join=round, line cap=round, thick, blue]
				\def\a{3}
				\path 	(0:0) coordinate (A)
				++(0:\a) coordinate (D)
				++(-130:\a/2) coordinate (C)
				($(A)+(C)-(D)$) coordinate (B)
				($(A)+(90:\a)$) coordinate (A')
				($(B)+(90:\a)$) coordinate (B')
				($(C)+(90:\a)$) coordinate (C')
				($(D)+(90:\a)$) coordinate (D');
				\draw[dashed] (B)--(A)--(D) (A)--(A') (A)--(C);
				\draw[] (C)--(C') (D)--(D') (B)--(B') (B)--(C)--(D) (A')--(B')--(C')--(D')--cycle ;
				\foreach \x/\y in {A/180,B/180,C/0,D/0,A'/180,B'/180,C'/0,D'/0}
				\fill[black] (\x) circle (1pt) ($(\y:3mm)+(\x)$) node {$\x$};	
			\end{tikzpicture}
		\end{center}
		Vì $ABCD.A'B'C'D'$ là hình lập phương nên ta có $\overrightarrow{B'C'} = \overrightarrow{AD}$.\\
		Do đó $\left(\overrightarrow{AC},\overrightarrow{B'C'} \right) = \left(\overrightarrow{AC},\overrightarrow{AD} \right) = \widehat{CAD}$.\\
		Tam giác $ADC$ vuông cân tại $D$ nên $\widehat{CAD} = 45^\circ$.\\
		Vậy góc giữa hai vec-tơ $\overrightarrow{AC}$ và $\overrightarrow{B'C'}$ bằng $45^\circ$.
	}
\end{ex}
%%%%%----------Câu 19
\begin{ex}%[2H2H1-2]
	Cho hình lăng trụ $ABC.A'B'C'$. Gọi $O$ là giao điểm của $AB'$ và $A'B$, $M$ là trung điểm của $AB$. Biết $\overrightarrow{CC'}=k\overrightarrow{OM}$. Giá trị $k$ bằng
	\shortans{$-2$}
	\loigiai{
		\begin{center}
			\begin{tikzpicture}[line cap=round, line join=round, font=\footnotesize, thick, blue, >=stealth]
				\def\a{3} \def\h{3.5}
				\path 	(0:0) coordinate (A)
				++(0:\a) coordinate (C)
				++(-150:3*\a/4) coordinate (B)
				($(A)+(80:\h)$) coordinate (A')
				($(B)+(80:\h)$) coordinate (B')
				($(C)+(80:\h)$) coordinate (C')
				(intersection of A--B' and B--A') coordinate (O)
				(barycentric cs:A=1,B=1) coordinate (M)
				;
				\draw[dashed] (A)--(C);
				\draw[] (C)--(C') (B)--(B') (A)--(A') (A)--(B)--(C) (A)--(B') (B)--(A') (A')--(B')--(C')--cycle;
				\draw[ultra thick,magenta,->] (O)--(M);
				\draw[ultra thick,magenta,->] (C)--(C');
				\foreach \x/\y in {A/180,B/-45,C/0,A'/180,B'/-45,C'/0, O/0,M/-100}
				\fill[black] (\x) circle (1pt) ($(\y:3mm)+(\x)$) node {$\x$};	
			\end{tikzpicture}
		\end{center}
		Ta có tứ giác $ABB'A'$ là hình bình hành nên $O$ là trung điểm của $AB'$ và $M$ là trung điểm của $AB$.\\ 
		Do đó $OM$ là đường trung bình của tam giác $AB'B$.\\
		Suy ra $BB'\parallel OM$ và $BB' = 2OM$.\\
		Tứ giác $BCC'B'$ là hình bình hành nên $BB'\parallel CC'$ và $BB'=CC'$. Do đó $CC'=2OM$.\\
		Mặt khác $\overrightarrow{CC'}$ và $\overrightarrow{OM}$ ngược hướng nên $\overrightarrow{CC'}=\left(-2\right)\overrightarrow{OM}$.\\ Vậy $k=-2$.
	}
\end{ex}
%%%%%----------Câu 20
\begin{ex}%[2H2H1-4]
	\immini{
		Ba lực  $\overrightarrow{F}_1$, $\overrightarrow{F}_2$, $\overrightarrow{F}_3$ cùng tác động vào một vật có phương đôi một vuông góc và có độ lớn lần lượt là $2\rm\,N$; $3\rm\,N$; $4\rm\,N$. Hợp lực của ba lực đã cho có độ lớn bao nhiêu Niu-tơn ({\it kết quả làm tròn đến một chữ số thập phân})?
	}{
		\begin{tikzpicture}[scale=1,>=stealth, font=\footnotesize, line join=round, line cap=round, thick]
			\path 
			(0:0) coordinate (O)
			($(O)+(-135:2)$) coordinate (A)
			($(O)+(0:3)$) coordinate (B)
			($(O)+(90:2.5)$) coordinate (C)
			;
			\draw[->] (O)--(A) node [pos=.4,left]{$\overrightarrow{F}_1 $};
			\draw[->] (O)--(B)node [pos=.5,above]{$\overrightarrow{F}_2 $};
			\draw[->] (O)--(C)node [pos=.5,left]{$\overrightarrow{F}_3 $};
			\draw pic[draw,angle radius=2mm]{right angle=A--O--B};
			\draw pic[draw,angle radius=2mm]{right angle=B--O--C};
			\draw pic[draw,angle radius=2mm]{right angle=C--O--A};
		\end{tikzpicture}
	}
	\shortans{$5{,}4$}
	\loigiai{
		\immini{
			Vẽ $\overrightarrow{OA}= \overrightarrow{F}_1$, $\overrightarrow{OB}= \overrightarrow{F}_2$, $\overrightarrow{OC}= \overrightarrow{F}_3$.\\
			Dựng hình chữ nhật $OABD$ và hình chữ nhật $ODEC$.\\
			Hợp lực tác động lên vật 
			$\overrightarrow{F}= \overrightarrow{F}_1+\overrightarrow{F}_2+\overrightarrow{F}_3 =
			\overrightarrow{OA}+\overrightarrow{OB}+\overrightarrow{OC}$.\\
			Theo quy tắc hình hộp ta có:\\
			$\left|\overrightarrow{F} \right|=
			\left|\overrightarrow{OA}+\overrightarrow{OB}+\overrightarrow{OC} \right|=
			\left|\overrightarrow{OE} \right| =OE$.\\
			Ta lại có $OE^2 = OA^2 + OB^2 + OC^2 = 29 $.\\
			$\Rightarrow OE =\sqrt{29}\approx 5{,}4$.\\
			Vậy hợp lực của ba lực đã cho có độ lớn bằng $5{,}4$ N. 
		}{
			
			\begin{tikzpicture}[scale=1,>=stealth, font=\footnotesize, line join=round, line cap=round, thick]
				\path 
				(0:0) coordinate (O)
				($(O)+(-135:2)$) coordinate (A)
				($(O)+(0:3)$) coordinate (B)
				($(O)+(90:2.5)$) coordinate (C)
				($(A)+(B)-(O)$) coordinate (D)
				($(C)+(D)-(O)$) coordinate (E)
				(intersection of O--B and D--E) coordinate (M)
				;
				\draw[->] (O)--(A) node [pos=.4,left]{$\overrightarrow{F}_1 $};
				\draw[->] (M)--(B)node [pos=.3,above]{$\overrightarrow{F}_2 $};
				\draw[->] (O)--(C)node [pos=.5,left]{$\overrightarrow{F}_3 $};
				\draw[->] (O)--(E)node [pos=.6,left]{$\overrightarrow{F} $};
				\draw (A)--(D)--(B) (C)--(E)--(D) (O)--(D);
				\draw[dashed] (O)--(M);
				\foreach \x/\y in {A/90,B/60,C/30,D/-30,E/30,O/135}
				\fill[black] (\x) circle (1pt) ($(\y:3mm)+(\x)$) node {$\x$};
			\end{tikzpicture}
		}
		
	}
\end{ex}
%%%%%----------Câu 21
\begin{ex}%[2H2H1-3]
	Trong không gian, cho hai vec-tơ $\overrightarrow{a}$ và $\overrightarrow{b}$. Biết  $\left|\overrightarrow{a} \right| = 1$, $\left|\overrightarrow{b} \right| = 2$ và góc giữa hai vec-tơ đó bằng $60^\circ$. Tính $\left(3\overrightarrow{a} + 2\overrightarrow{b} \right)^2$.
	\shortans{$37$}
	\loigiai{
		\begin{eqnarray*}
			\left(3\overrightarrow{a} + 2\overrightarrow{b} \right)^2
			&=& \left(3\overrightarrow{a} \right)^2 + 2\cdot 3\overrightarrow{a} \cdot 2\overrightarrow{b} + \left(2\overrightarrow{b} \right)^2\\
			&=& 9 \left|\overrightarrow{a} \right|^2 + 12\cdot\overrightarrow{a} \cdot \overrightarrow{b}+ 4\left|\overrightarrow{b} \right|^2\\
			&=& 9 + 12 \left|\overrightarrow{a} \right| \cdot\left|\overrightarrow{b} \right|\cdot \cos\left(\overrightarrow{a},\overrightarrow{b} \right)+ 16\\
			&=& 9 + 12\cdot 1\cdot 2\cdot \cos 60^\circ + 16\\
			&=& 9 + 12 + 16 = 37.
		\end{eqnarray*}
	}
\end{ex}
%%%%%----------Câu 22
\begin{ex}%[2H2H1-3]
	Cho tứ diện đều $ABCD$ có cạnh bằng $6$ và $M$ là trung điểm của $CD$. Tính tích vô hướng $\overrightarrow{AB}\cdot \overrightarrow{AM}$.
	\shortans{$18$}
	\loigiai{
		\immini{
			Vì $M$ là trung điểm $CD$ nên ta có $\overrightarrow{AC} + \overrightarrow{AD} = 2\overrightarrow{AM}$.
			\begin{eqnarray*}
				\text{Ta có: } \overrightarrow{AB}\cdot \overrightarrow{AC} 
				&=& \left|\overrightarrow{AB}\right| \cdot \left|\overrightarrow{AC}\right| \cdot \cos \left(\overrightarrow{AB},\overrightarrow{AC}\right)\\
				&= & 6\cdot 6 \cdot \cos 60^\circ = 18. 
			\end{eqnarray*}
			Tương tự $\overrightarrow{AB}\cdot \overrightarrow{AD} = 18$.
			\begin{eqnarray*}
				\text{Suy ra }\overrightarrow{AB}\cdot \overrightarrow{AM} 
				&=& \overrightarrow{AB}\cdot \dfrac{1}{2}\left(\overrightarrow{AC} + \overrightarrow{AD}\right)\\
				&=& \dfrac{1}{2}\left(\overrightarrow{AB}\cdot\overrightarrow{AC}+ \overrightarrow{AB}\cdot\overrightarrow{AD}\right)\\
				&=& \dfrac{1}{2} \left(18 + 18\right)= 18.
			\end{eqnarray*}
			
		}{	
			\begin{tikzpicture}[scale=1,>=stealth, font=\footnotesize, line join=round, line cap=round, thick]
				\def\a{1} \def\h{4}
				\path 	(0:0) coordinate (B)
				(0:4*\a) coordinate (D)
				(320:2*\a) coordinate (C)
				($(C)!0.5!(D)$) coordinate (M)
				($(B)!2/3!(M)$) coordinate (H)
				($(H)+(90:\h)$) coordinate (A);
				\draw[ ] (A)--(B)--(C)--(D)--cycle (A)--(C) (A)--(M);
				\draw[dashed,thick] (B)--(D);
				\foreach \x / \y in {A/120,B/150,C/-135,M/-60,D/-30}
				\fill (\x) circle (1pt) ($(\x)+(\y:3mm)$) node {$\x$};
			\end{tikzpicture}
		}
	}
\end{ex}
\Closesolutionfile{ans}
\inputansbox[3]{6}{ans/ans-2-B6-KQ}
